\section*{अध्याय \ref{ch:b2:hermitian} --- हर्मीशियन रूप}

\begin{solution}{exo:b2:hermitian:1}
$\langle x, y\rangle = \conj{1}\cdot\iu + \conj{\iu}\cdot 1 = \iu -
\iu = 0$: लांबिक। $\norm x = \norm y = \sqrt{1 + 1} = \sqrt2$। प्रसामान्य लांबिक आधार: $\bigl(\frac{1}{\sqrt2}(1, \iu),\;
\frac{1}{\sqrt2}(\iu, 1)\bigr)$ --- अर्थात् सामान्यीकृत
युग्म स्वयं।
\end{solution}

\begin{solution}{exo:b2:hermitian:2}
पहला: अपने संयुग्मी \omterm{def:b2:linalg:transpose}{परिवर्त} के बराबर (वास्तविक विकर्ण, $\conj{\iu} =
-\iu$ आपस में बदले
हुए): अतः \omterm{def:b2:hermitian:adjoint}{हर्मीशियन} (इसलिए प्रसामान्य); पर \omterm{def:b2:hermitian:adjoint}{एकात्मक} नहीं ($A^\dagger A
\neq I$: स्तंभ इकाई
नहीं)।

दूसरा: $A^\dagger A = \frac12\begin{pmatrix} 1 & -\iu\\ -\iu &
1\end{pmatrix}\begin{pmatrix} 1 & \iu\\ \iu & 1\end{pmatrix} =
\frac12\begin{pmatrix} 2 & 0\\ 0 & 2\end{pmatrix} = I$: अतः \omterm{def:b2:hermitian:adjoint}{एकात्मक} (इसलिए प्रसामान्य); \omterm{def:b2:hermitian:adjoint}{हर्मीशियन} नहीं।

तीसरा: $A^\dagger A = E_{22} \neq E_{11} = AA^\dagger$: अतः प्रसामान्य नहीं (इसलिए न \omterm{def:b2:hermitian:adjoint}{हर्मीशियन} न \omterm{def:b2:hermitian:adjoint}{एकात्मक}) --- यही
मानक शून्यंभावी प्रतिउदाहरण है।
\end{solution}

\begin{solution}{exo:b2:hermitian:3}
अद्वितीयता: $H^\dagger = H$, $K^\dagger = K$ के साथ $A = H + \iu K$ $A^\dagger = H - \iu K$ को बाध्य कर देता है, अतः $H = \frac{A + A^\dagger}{2}$, $K = \frac{A - A^\dagger}{2\iu}$;
और ये सूत्र \omterm{def:b2:hermitian:adjoint}{हर्मीशियन} हैं (जाँच: $\bigl(\frac{A - A^\dagger}{2\iu}\bigr)^\dagger = \frac{
A^\dagger - A}{-2\iu} = K$) तथा $A$ को पुनः बना देते हैं:
अर्थात् अस्तित्व।

प्रसामान्यता: $A^\dagger A - AA^\dagger = (H - \iu K)(H + \iu K) - (H
+ \iu K)(H - \iu K) = 2\iu(HK - KH)$: यह लुप्त होता है तभी जब $HK = KH$।
\end{solution}

\begin{solution}{exo:b2:hermitian:4}
$\chi_A = (X-2)^2 - 1$: \omterm{def:b2:reduction:eigen}{अभिलक्षणिक मान} $3$ और $1$। \omterm{def:b2:reduction:eigen}{अभिलक्षणिक सदिश}, सीधे परिकलन से:
\[
A\begin{pmatrix}1\\ -\iu\end{pmatrix}
= \begin{pmatrix} 2 + \iu(-\iu)\\ -\iu + 2(-\iu)\end{pmatrix}
= \begin{pmatrix} 3\\ -3\iu \end{pmatrix}
= 3\begin{pmatrix}1\\ -\iu\end{pmatrix},
\qquad
A\begin{pmatrix}1\\ \iu\end{pmatrix}
= \begin{pmatrix} 2 + \iu\cdot\iu\\ -\iu + 2\iu\end{pmatrix}
= \begin{pmatrix}1\\ \iu\end{pmatrix}.
\]
प्रसामान्य लांबिक अभिलक्षणिक आधार: $u_1 = \frac{1}{\sqrt2}(1, -\iu)$ (\omterm{def:b2:reduction:eigen}{अभिलक्षणिक मान} $3$), $u_2 = \frac{1}{\sqrt2}(1, \iu)$
(\omterm{def:b2:reduction:eigen}{अभिलक्षणिक मान} $1$); और लांबिकता \cref{exo:b2:hermitian:1} की तरह। वर्णक्रमीय प्रक्षेपों $A^k = 3^k u_1u_1^\dagger + 1^k\,
u_2u_2^\dagger$
के द्वारा घातें:
\[
A^k = U\begin{pmatrix} 3^k & 0\\ 0 & 1\end{pmatrix}U^\dagger
= \frac{3^k}{2}\begin{pmatrix} 1 & \iu\\ -\iu & 1\end{pmatrix}
+ \frac{1}{2}\begin{pmatrix} 1 & -\iu\\ \iu & 1\end{pmatrix}
= \frac12\begin{pmatrix}
3^k + 1 & (3^k - 1)\iu\\
-(3^k-1)\iu & 3^k + 1
\end{pmatrix}.
\]
($k = 1$ जाँचिए: $A$ पुनः मिल जाता है।)
\end{solution}

\begin{solution}{exo:b2:hermitian:5}
\omterm{def:b2:metric:compact}{संहत}: $\leq 1$ में संवृत (\omterm{def:b2:metric:continuity}{संतत} $U \mapsto
U^\dagger U$ के अंतर्गत $I$ का पूर्वप्रतिबिंब) और
परिबद्ध (स्तंभ इकाई सदिश हैं: प्रविष्टियों का मापांक $\mathcal{M}_n(\C) \simeq \R^{2n^2}$)।

\omterm{def:b2:reduction:eigen}{अभिलक्षणिक मान}: किसी भी $\abs\lambda = 1$ के लिए $\operatorname{diag}(\lambda, 1, \dots, 1)$ \omterm{def:b2:hermitian:adjoint}{एकात्मक} है। \omterm{def:b2:linalg:det}{सारणिक}: $\abs{\det U}^2 = \det
U^\dagger \det U = \det(U^\dagger U) = 1$ ($\det A^\dagger =
\conj{\det A}$ का
उपयोग करते हुए): अतः $\det U$ इकाई वृत्त पर पड़ता है।
\end{solution}

\begin{solution}{exo:b2:hermitian:6}
वर्णक्रमीय प्रमेय से $H$ के किसी प्रसामान्य लांबिक अभिलक्षणिक आधार में
काम कीजिए: सब कुछ अदिशों $\lambda \in \R$ (अर्थात् अभिलक्षणिक मानों) तक सिमट जाता है।
$H + \iu I$ के \omterm{def:b2:reduction:eigen}{अभिलक्षणिक मान} $\lambda + \iu \neq 0$ हैं: अतः प्रतिलोमनीय। $U$ के \omterm{def:b2:reduction:eigen}{अभिलक्षणिक मान} $\mu = \frac{\lambda - \iu}{\lambda + \iu}$
हैं, जिनका मापांक $1$ है (वास्तविक $\lambda$ के लिए $\abs{\lambda - \iu} = \abs{\lambda + \iu}$): और $U^\dagger U = I$ सत्य है
क्योंकि $U$ इकाई-मापांक अभिलक्षणिक मानों के साथ \omterm{def:b2:hermitian:adjoint}{एकात्मक} रूप से \omterm{def:b2:reduction:diag}{विकर्णनीय}
है (वह चुने हुए प्रसामान्य लांबिक आधार में विकर्ण है)। और $\mu = 1$ $-\iu = \iu$ को
बाध्य कर देता, जो असंभव है, अतः $1 \notin \operatorname{Sp} U$। (केली रूपांतरण \omterm{def:b2:hermitian:adjoint}{हर्मीशियन} को
एकात्मक-में-से-एक-बिंदु पर भेजता है --- अर्थात् $\R$ से वृत्त पर जाने वाले
प्रतिचित्रण का आव्यूह-रूप।)
\end{solution}

\begin{solution}{exo:b2:hermitian:7}
किसी अभिलक्षणिक युग्म $Ax = \lambda x$ ($x \neq 0$) के लिए: $\lambda\norm x^2 =
\langle x, Ax\rangle > 0$, अतः $\lambda > 0$ (जो पहले से ही
वास्तविक है, \cref{prop:b2:hermitian:eigenvalues})। वर्गमूल: किसी वर्णक्रमीय आधार में $B = U\operatorname{diag}(\sqrt{\lambda_i})U^\dagger$: यह \omterm{def:b2:hermitian:adjoint}{हर्मीशियन},
धनात्मक निश्चित है और $B^2 = A$। \omterm{def:b2:linalg:det}{सारणिक}: धनात्मक अभिलक्षणिक मानों का गुणनफल।
\end{solution}

\begin{solution}{exo:b2:hermitian:8}
\begin{enumerate}
  \item $\norm{u(x)}^2 = \langle u(x), u(x)\rangle = \langle x,
        u^*u(x)\rangle = \langle x, uu^*(x)\rangle =
        \norm{u^*(x)}^2$। इसे प्रसामान्य $u -
        \lambda\,\mathrm{id}$ पर लगाइए (उसका \omterm{def:b2:hermitian:adjoint}{सहसंलग्न} $u^* -
        \conj\lambda$ है, और
        प्रसामान्यता विरासत में मिलती है): $\norm{(u -
        \lambda)x} = \norm{(u^* - \conj\lambda)x}$, अतः दोनों केंद्रक मेल
        खाते हैं।
  \item अभिलक्षणिक सदिशों $u(x) = \lambda x$, $u(y) = \mu y$ ($\lambda \neq \mu$) के लिए: (1) का उपयोग करके $u^*(x) = \conj\lambda x$;
        फिर
        \[
        \lambda\langle y, x\rangle = \langle y, u(x)\rangle
        = \langle u^*(y), x\rangle = \langle \conj\mu\, y, x\rangle
        = \mu \langle y, x\rangle ,
        \]
        अतः $\langle y, x\rangle = 0$। $E_\lambda^\perp$ का स्थायित्व: $x \perp E_\lambda$ और $z \in E_\lambda$ के लिए $\langle z,
        u(x)\rangle = \langle u^*(z), x\rangle =
        \langle\conj\lambda z, x\rangle = 0$।
  \item विमा पर आगमन: $\C$ पर $u$ का कोई \omterm{def:b2:reduction:eigen}{अभिलक्षणिक सदिश} $e_1$ है (उसे
        सामान्यीकृत कीजिए); उसका लांबिक पूरक $u$ के अंतर्गत स्थायी है
        ((2) से) \emph{और} $u^*$ के अंतर्गत भी (भूमिकाएँ बदलकर वही तर्क),
        अतः प्रतिबंधन प्रसामान्य है: आगमन कीजिए और प्रसामान्य लांबिक
        अभिलक्षणिक आधार जोड़ते जाइए। विलोमतः \omterm{def:b2:hermitian:adjoint}{एकात्मक} रूप से \omterm{def:b2:reduction:diag}{विकर्णनीय} $u = UDU^\dagger$
        $u^*u
        = U\conj D D U^\dagger = U D\conj D U^\dagger = uu^*$ पूरा करता है: अर्थात् प्रसामान्य।
\end{enumerate}
\end{solution}

\begin{solution}{exo:b2:hermitian:9}
\omterm{def:b2:reduction:eigen}{अभिलक्षणिक मान}: $u(x) = \lambda x$, $x \neq 0$ के लिए
\[
\lambda\norm x^2 = \langle x, u(x)\rangle
= \langle u^*(x), x\rangle = -\langle u(x), x\rangle
= -\conj{\langle x, u(x)\rangle} = -\conj\lambda\,\norm x^2 ,
\]
अतः $\lambda = -\conj\lambda$: अर्थात् विशुद्ध काल्पनिक। और चूँकि $(\iu
u)^* = -\iu\,u^*$ (\omterm{def:b2:hermitian:adjoint}{सहसंलग्न} अदिशों में
संयुग्मी-रैखिक है), $u^* = u$ $(\iu u)^* = -\iu u$ देता है: अतः प्रतिचित्रण $u \mapsto \iu u$ \omterm{def:b2:hermitian:adjoint}{हर्मीशियन} को
विषम-हर्मीशियन पर भेजता है और उसका प्रतिलोम $w \mapsto -\iu
w$ है: यानी एकैकी आच्छादक।
और कोई विषम-हर्मीशियन $u$ $u^*u = -u^2 =
uu^*$ पूरा करता है: अर्थात् प्रसामान्य, इसलिए
\cref{exo:b2:hermitian:8} से \omterm{def:b2:hermitian:adjoint}{एकात्मक} रूप से \omterm{def:b2:reduction:diag}{विकर्णनीय}।
\end{solution}

\begin{solution}{exo:b2:hermitian:10}
\begin{enumerate}
  \item $S$ किसी प्रसामान्य लांबिक आधार का क्रमचय करता है: $\norm{Sx} = \norm
        x$, अतः $S$
        \omterm{def:b2:hermitian:adjoint}{एकात्मक} है। निर्देशांकों को $n$ के सापेक्ष $j = 0,
        \dots, n-1$ से अंकित कीजिए:
        $(Sx)_j = x_{j-1}$, अतः $(f_k)_j = \frac{\omega^{jk}}{\sqrt n}$ के लिए:
        \[
        (Sf_k)_j = \frac{\omega^{(j-1)k}}{\sqrt n}
        = \omega^{-k}\,(f_k)_j :
        \qquad Sf_k = \omega^{-k}f_k .
        \]
        प्रसामान्य लांबिकता: $\langle f_k, f_l\rangle = \frac1n
        \sum_j \omega^{j(l-k)} = \delta_{kl}$ (किसी अतुच्छ इकाई-मूल का गुणोत्तर योग
        लुप्त हो जाता है)।
  \item $Cf_k = \sum_j c_j S^jf_k = \bigl(\sum_j
        c_j\omega^{-jk}\bigr)f_k = \widehat c(k)\,f_k$: अर्थात् हर \omterm{def:b2:structures:generated}{चक्रीय} आव्यूह प्रसामान्य लांबिक फूरिये आधार में
        विकर्ण है, इसलिए प्रसामान्य, और उसका \omterm{def:b2:reduction:eigen}{वर्णक्रम} $\{\widehat c(k)\}$ है।
        (आधार-परिवर्तन विविक्त फूरिये रूपांतरण है: संवलन गुणन बन जाता
        है।)
\end{enumerate}
\end{solution}

\begin{solution}{exo:b2:hermitian:11}
($\Leftarrow$) मान लीजिए $P^2 = P = P^*$। हर $v$ $Pv \in \operatorname{im} P$ तथा $P(v - Pv) =
0$ के साथ $v = Pv
+ (v - Pv)$ में बँट जाता है। दोनों
टुकड़े लांबिक हैं: किसी भी $x, y$ के लिए
\[
\langle Px, (I - P)y\rangle = \langle x, P(I-P)y\rangle
= \langle x, (P - P^2)y\rangle = 0 :
\]
$\ker P \perp \operatorname{im} P$, अतः $P$ अपने प्रतिबिंब पर लांबिक प्रक्षेप है। ($\Rightarrow$) यदि $P$ $F = \operatorname{im}P$
पर लांबिक प्रक्षेप हो: तो सभी $x, y$ के लिए $\langle Px, y\rangle = \langle Px, Py\rangle$ (घटक $y - Py \perp F$ हट जाता है) और
सममित रूप से $\langle
x, Py\rangle = \langle Px, Py\rangle$: $\langle Px, y\rangle =
\langle x, Py\rangle$, अर्थात् $P^* = P$। आव्यूह: $\C v$ ($\norm v = 1$) पर: $Px = v\,\langle v, x\rangle$,
अर्थात् $P =
vv^\dagger$; और प्रसामान्य लांबिक $\operatorname{Vect}(v_1, \dots, v_k)$ वाली उपसमष्टि पर: $P = \sum_i v_iv_i^\dagger$।
\end{solution}

\begin{solution}{exo:b2:hermitian:12}
$A = U D U^\dagger$ का विकर्णन कीजिए (वर्णक्रमीय प्रमेय), जहाँ $D$ विकर्ण है और उसकी
प्रविष्टियाँ $\lambda_i$ में से हैं। तब $P_i = L_i(A) = U
L_i(D)U^\dagger$, और $L_i(D)$ विकर्ण है जिसकी प्रविष्टियाँ
$L_i(\lambda_j) = \delta_{ij}$ हैं: अर्थात् ठीक $E_i$ के खाँचों पर एक। अतः $P_i$ \omterm{def:b2:hermitian:adjoint}{हर्मीशियन} ($L_i$ वास्तविक,
$D$ वास्तविक), वर्गसम है, जिसका प्रतिबिंब $E_i$ और केंद्रक $\bigoplus_{j\neq i}E_j = E_i^\perp$ है
(अभिलक्षणिक समष्टियों की लांबिकता): यानी $E_i$ पर लांबिक प्रक्षेप (\cref{exo:b2:hermitian:11})।
असंयुक्त विकर्ण-प्रतिरूप $P_iP_j = 0$ देते हैं ($i \neq j$); $\sum_i L_i = 1$ (घात $< p$, और $p$
बिंदुओं पर मान $1$), अतः $\sum P_i = I$; और $\sum_i
\lambda_iL_i(\lambda_j) = \lambda_j$ $A = \sum
\lambda_iP_i$ देता है। किसी भी बहुपद $f$
के लिए: $f(D)$ का विकर्ण $f(\lambda_j)$ है, अतः
\[
f(A) = \sum_{i=1}^{p} f(\lambda_i)\,P_i :
\]
अर्थात् $A$ के फलन वर्णक्रमीय ढंग से परिकलित होते हैं --- और तृतीय वर्ष
का खंड इसी कलन को \omterm{def:b2:metric:continuity}{संतत} $f$ तथा उससे आगे तक बढ़ा देता है।
\end{solution}

\begin{solution}{pb:b2:hermitian:1}
\textbf{1.} $\conj{\langle x, Ax\rangle} = \langle Ax, x\rangle
= \langle x, A^*x\rangle = \langle x, Ax\rangle$: अर्थात् वास्तविक। $x = \sum c_ie_i$ लिखने पर:
\[
R_A(x) = \frac{\sum_i\lambda_i\abs{c_i}^2}
{\sum_i\abs{c_i}^2} ,
\]
जो अभिलक्षणिक मानों का भारित औसत है: अतः वह $\intcc{\lambda_n}{\lambda_1}$ में पड़ता है, और परिबंध
$e_1$ तथा $e_n$ पर प्राप्त होते हैं।

\textbf{2.} वास्तविक $t$ और किसी भी $v$ के लिए $R_A(x + tv) =
\frac{N(t)}{D(t)}$ का प्रसार कीजिए
\[
N(t) = \langle x, Ax\rangle + 2t\Re\langle v, Ax\rangle +
t^2\langle v, Av\rangle,
\quad
D(t) = \norm x^2 + 2t\Re\langle v, x\rangle + t^2\norm v^2 .
\]
$t = 0$ पर अवकलज
\[
\frac{2}{\norm x^2}\,
\Re\bigl\langle v,\ Ax - R_A(x)\,x\bigr\rangle .
\]
है। वह सभी $v$ के लिए लुप्त होता है तभी जब सभी $v$ के लिए $\Re\langle v, w\rangle = 0$, जहाँ $w = Ax - R_A(x)x$;
और $v$ के स्थान पर $\iu v$ रखने से काल्पनिक भाग भी मर जाता है: $w = 0$,
अर्थात् $Ax = R_A(x)x$। अतः $R_A$ के क्रांतिक बिंदु ठीक \omterm{def:b2:reduction:eigen}{अभिलक्षणिक सदिश} हैं, और
क्रांतिक मान वही \omterm{def:b2:reduction:eigen}{अभिलक्षणिक मान}।

\textbf{3.} $x = \sum_{i\leq k}c_ie_i \in V_k$ के लिए: $R_A(x)$ $\lambda_1, \dots, \lambda_k$ का भारित औसत है, अतः $\geq \lambda_k$, और $e_k$ पर समता:
$\min_{V_k} R_A =
\lambda_k$। सममित रूप से $W_k$ पर औसत में $\lambda_k, \dots, \lambda_n$ आते हैं: $\max_{W_k}R_A = \lambda_k$।

\textbf{4.} मान लीजिए $\dim V = k$। तब $\dim V + \dim W_k = n + 1 >
n$, अतः कोई इकाई $x \in V \cap W_k$ है और $R_A(x) \leq
\lambda_k$ (प्रश्न 3):
इसलिए ऐसे प्रत्येक $V$ के लिए $\min_{V}R_A \leq \lambda_k$। और चूँकि $V_k$ $\lambda_k$ प्राप्त कर लेता है,
महत्तम-न्यूनतम $\lambda_k$ के बराबर है। न्यूनतम-महत्तम सूत्र वही तर्क है, बस
भूमिकाएँ बदली हुई ($\dim W = n - k + 1$ $W \cap V_k \neq
\{0\}$ को बाध्य करता है, अतः $\max_W R_A \geq \lambda_k$, जो $W_k$ पर
प्राप्त होता है)।

\textbf{5.} \omterm{def:b2:funcseq:def}{बिंदुवार} $R_B(x) = R_A(x) + \frac{\langle x,
(B-A)x\rangle}{\norm x^2} \geq R_A(x)$। किसी भी $k$-विमीय $V$ पर $\min$ लेकर और फिर $V$ पर
$\max$ लेकर: प्रश्न 4 से $\lambda_k(B) \geq \lambda_k(A)$।

\textbf{6.} ग्रासमान: $\dim(V\cap W) = \dim V + \dim W -
\dim(V + W) \geq \dim V + \dim W - n > 0$। इसे दो बार लगाने पर:
\[
\dim(V_1\cap V_2\cap V_3)
\geq \dim(V_1\cap V_2) + \dim V_3 - n
\geq \dim V_1 + \dim V_2 + \dim V_3 - 2n .
\]

\textbf{7.} उपसमष्टियों $W_i(A)$, $W_j(B)$ (प्रश्न 3, $A$ तथा $B$ के लिए) और $V_{i+j-1}(A+B)$ की
विमाएँ $(n-i+1) +
(n-j+1) + (i+j-1) = 2n + 1 > 2n$ हैं: अतः प्रश्न 6 से तीनों में कोई इकाई सदिश $x$ है। तब
\[
\lambda_{i+j-1}(A+B) \leq R_{A+B}(x)
= R_A(x) + R_B(x) \leq \lambda_i(A) + \lambda_j(B),
\]
जिसमें बाईं असमिका $x \in V_{i+j-1}(A+B)$ के कारण (प्रश्न 3) और दाईं दोनों $W$ के कारण है।

\textbf{8.} $E$ के किसी वर्णक्रमीय आधार में: $\norm{Ex}^2 = \sum
\lambda_k(E)^2\abs{c_k}^2 \leq \max_k\lambda_k(E)^2\,\norm
x^2$, जो संगत \omterm{def:b2:reduction:eigen}{अभिलक्षणिक सदिश}
पर प्राप्त होता है: $\vertiii E_2
= \max_k\abs{\lambda_k(E)}$। $j = 1$ के साथ वाइल: $\lambda_k(A+E)
\leq \lambda_k(A) + \lambda_1(E) \leq \lambda_k(A) + \vertiii
E_2$; और इसे $(A+E) + (-E)$ पर लगाने पर:
$\lambda_k(A) \leq
\lambda_k(A+E) + \vertiii E_2$। दोनों मिलाकर: $\abs{\lambda_k(A+E) - \lambda_k(A)} \leq \vertiii E_2$।

\textbf{9.} निचले परिबंध: $P \geq 0$ और प्रश्न 5। ऊपरी: $P$ की कोटि $1$ है, अतः
$\lambda_2(P) = 0$; और $i = k-1$, $j =
2$ के साथ वाइल:
\[
\lambda_k(A + P) \leq \lambda_{k-1}(A) + \lambda_2(P)
= \lambda_{k-1}(A) .
\]

\textbf{10.} $A + E = \begin{pmatrix} 2 & 1+\iu\\ 1-\iu &
2\end{pmatrix}$: \omterm{def:b2:reduction:charpoly}{अभिलक्षणिक बहुपद} $(2-\lambda)^2 -
\abs{1+\iu}^2 = (2-\lambda)^2 - 2$, \omterm{def:b2:reduction:eigen}{वर्णक्रम} $\{2 + \sqrt2,\ 2
- \sqrt2\}$। $\operatorname{Sp}A = \{3, 1\}$ के विरुद्ध:
\[
\abs{(2+\sqrt2) - 3} = \abs{(2-\sqrt2) - 1} = \sqrt2 - 1
\approx 0.414 \leq 1 = \vertiii E_2 . \checkmark
\]

\textbf{11.} $\C^{n-1} = \operatorname{Vect}(e_1, \dots,
e_{n-1})$ को $\C^n$ के भीतर देखिए (मानक आधार): वहाँ $x$ के लिए $\langle x, Bx\rangle = \langle x, Ax\rangle$,
अतः $R_B$ $R_A$ का प्रतिबंधन है। ऊपरी परिबंध: $\lambda_k
(B)$ के लिए महत्तम-न्यूनतम $\C^{n-1}$
\emph{की} $k$-विमीय उपसमष्टियों पर घूमता है, जो $\C^n$ की उपसमष्टियों का
उपकुल है: अतः $\lambda_k(B) \leq
\lambda_k(A)$। निचला परिबंध: $\lambda_k(B)$ के लिए न्यूनतम-महत्तम $\C^{n-1}$ की विमा
$(n-1)-k+1 =
n-k$ वाली उपसमष्टियों पर घूमता है; और हर एक $\C^n$ की विमा $n -
(k+1) + 1$ वाली
उपसमष्टि भी है, अतः उसका महत्तम $\geq \lambda_{k+1}(A)$ है: $\lambda_k(B) \geq \lambda_{k+1}(A)$।

\textbf{12.} एक-एक करके पंक्तियाँ/स्तंभ हटाइए और प्रश्न 11 को शृंखलाबद्ध
कीजिए: हर हटाव निचला सूचकांक एक से खिसका देता है, जिससे $\lambda_{k+m}(A) \leq \lambda_k(B) \leq \lambda_k(A)$ मिलता है।

\textbf{13.} $A$ का किसी क्रमचय आव्यूह (जो \omterm{def:b2:hermitian:adjoint}{एकात्मक} है) से संयुग्मन न
\omterm{def:b2:reduction:eigen}{वर्णक्रम} बदलता है न विकर्ण प्रविष्टियों का बहुसमुच्चय: अतः मान लीजिए $d_1, \dots, d_k$
अग्र स्थानों पर हैं। तब अग्र $k\times k$ मुख्य उपआव्यूह $B$ के लिए $\operatorname{tr} B = \sum_{i\leq k}d_i$, और उसके
\omterm{def:b2:reduction:eigen}{अभिलक्षणिक मान} $\mu_i(B) \leq \lambda_i(A)$ पूरा करते हैं (प्रश्न 12): योग लेने पर $\sum_{i\leq k}d_i \leq \sum_{i\leq k}\lambda_i(A)$। और $k =
n$ पर
दोनों पक्ष $\operatorname{tr} A$ हैं।

\textbf{14.} $x_i = e_i$ लेने पर मान $\sum_{i\leq k}\lambda_i$ मिलता है: अतः महत्तम $\geq$ है। विलोमतः कोई
प्रसामान्य लांबिक कुल $(x_1, \dots, x_k)$ किसी प्रसामान्य लांबिक आधार तक बढ़ाया जा सकता
है, अर्थात् ऐसे \omterm{def:b2:hermitian:adjoint}{एकात्मक} $U$ तक जिसके पहले स्तंभ $x_i$ हों; तब $\sum_i\langle x_i, Ax_i\rangle$ $U^\dagger AU$ की
पहली $k$ विकर्ण प्रविष्टियों का योग है, जो प्रश्न 13 के अनुसार उसके $k$
सबसे बड़े अभिलक्षणिक मानों के योग से अधिक नहीं --- अर्थात् $\sum_{i\leq k}\lambda_i(A)$। इस प्रकार
काय फान का महत्तम-सिद्धांत निकल आता है।

\textbf{15.} अंतर्गुंथन ($\operatorname{Sp}A = \{2+\sqrt2, 2,
2-\sqrt2\}$, $\operatorname{Sp}B = \{3, 1\}$):
\[
2 \leq 3 \leq 2 + \sqrt2,
\qquad
2 - \sqrt2 \leq 1 \leq 2 . \checkmark
\]
विकर्ण $(2,2,2)$ के साथ शूर: $2 \leq 2+\sqrt2$; $4 \leq 4 +
\sqrt2$; $6 = 6$ (अनुरेख)। \checkmark

\textbf{16.} $b_{ii} - a_{ii} = \langle e_i, (B-A)e_i\rangle
\geq 0$; योग लेने पर अनुरेख मिल जाते हैं। और किसी भी $C$ के लिए:
$\langle x,
C^\dagger(B - A)Cx\rangle = \langle Cx, (B-A)(Cx)\rangle \geq
0$: $C^\dagger AC \leq C^\dagger BC$।

\textbf{17.} $A \geq 0$ स्पष्ट है; $B - A = \begin{pmatrix} 1 & 1\\
1 & 1\end{pmatrix}$ धनात्मक अर्धनिश्चित है (\omterm{def:b2:reduction:eigen}{अभिलक्षणिक मान}
$2,
0$): $A \leq B$। परंतु
\[
B^2 = \begin{pmatrix} 5 & 3\\ 3 & 2\end{pmatrix},
\qquad
B^2 - A^2 = \begin{pmatrix} 4 & 3\\ 3 & 2\end{pmatrix},
\qquad \det(B^2 - A^2) = -1 < 0 :
\]
धनात्मक अर्धनिश्चित नहीं। अतः वर्ग करना लोएव्नर क्रम का सम्मान नहीं करता।

\textbf{18.} मान लीजिए $S = \sqrt A$, $T = \sqrt B$ (\omterm{def:b2:hermitian:adjoint}{हर्मीशियन} धनात्मक अर्धनिश्चित, और \cref{exo:b2:hermitian:7} को
उसी वर्णक्रमीय सूत्र से अर्धनिश्चित तक बढ़ाया गया)। $T - S$ \omterm{def:b2:hermitian:adjoint}{हर्मीशियन} है; $\mu$
को कोई भी \omterm{def:b2:reduction:eigen}{अभिलक्षणिक मान} लीजिए और $v$ उसका इकाई \omterm{def:b2:reduction:eigen}{अभिलक्षणिक सदिश}। $T^2 - S^2 = T(T - S) + (T - S)S$ से:
\[
0 \leq \langle v, (B - A)v\rangle
= \langle Tv, (T-S)v\rangle + \langle (T-S)v, Sv\rangle
= \mu\bigl(\langle v, Tv\rangle + \langle v,
Sv\rangle\bigr)
\]
($\mu$ वास्तविक है)। यदि $\langle v, Tv\rangle + \langle v,
Sv\rangle > 0$ हो, तो $\mu \geq 0$। और यदि वह लुप्त हो, तो दोनों
अऋणात्मक पद लुप्त हैं; $\langle v, Tv\rangle =
\norm{T^{1/2}v}^2 = 0$ $Tv = 0$ को बाध्य करता है, वैसे ही $Sv = 0$, अतः $\mu v = (T - S)v = 0$
और $\mu = 0$। इस प्रकार $T -
S$ के सारे \omterm{def:b2:reduction:eigen}{अभिलक्षणिक मान} $\geq 0$ हैं: $\sqrt A \leq \sqrt B$।

\textbf{19.} $A^{-1/2}$ से \omterm{def:b2:quadratic:def}{सर्वांगसमता} (प्रश्न 16): $I \leq M
:= A^{-1/2}BA^{-1/2}$। अतः $M$ के सारे
\omterm{def:b2:reduction:eigen}{अभिलक्षणिक मान} $\geq 1$ हैं, इसलिए $M^{-1}$ के \omterm{def:b2:reduction:eigen}{अभिलक्षणिक मान} $\intoc{0}{1}$ में पड़ते हैं:
$M^{-1} \leq I$। परंतु $M^{-1} = A^{1/2}B^{-1}A^{1/2}$; और $A^{-1/2}$ से $M^{-1} \leq
I$ की \omterm{def:b2:quadratic:def}{सर्वांगसमता} $B^{-1} \leq A^{-1}$ देती है।

\textbf{20.} $\sqrt A^{-1}(AB)\sqrt A = \sqrt A\,B\sqrt A$: अतः $AB$ \omterm{def:b2:hermitian:adjoint}{हर्मीशियन} धनात्मक निश्चित $\sqrt
A\,B\sqrt A$ के सदृश है
(निश्चित: $\langle x, \sqrt AB\sqrt Ax\rangle =
\langle \sqrt Ax, B\sqrt Ax\rangle > 0$): इसलिए उसके \omterm{def:b2:reduction:eigen}{अभिलक्षणिक मान} वास्तविक और धनात्मक हैं। इसके
अतिरिक्त
\[
\langle x, \sqrt AB\sqrt Ax\rangle
\leq \lambda_{\max}(B)\,\norm{\sqrt Ax}^2
= \lambda_{\max}(B)\,\langle x, Ax\rangle
\leq \lambda_{\max}(A)\lambda_{\max}(B)\norm x^2 ,
\]
अतः $\lambda_{\max}(AB) = \max R_{\sqrt AB\sqrt A} \leq
\lambda_{\max}(A)\lambda_{\max}(B)$।

\textbf{21.} $v_k = (\sin j\theta_k)_j$ तथा सर्वसमिका $\sin((j-1)\theta) + \sin((j+1)\theta) =
2\sin(j\theta)\cos\theta$ के साथ: $2 \leq j \leq n-1$ के लिए
\[
(T_nv_k)_j = -\sin((j{-}1)\theta_k) + 2\sin(j\theta_k) -
\sin((j{+}1)\theta_k)
= (2 - 2\cos\theta_k)\sin(j\theta_k) .
\]
पंक्ति $1$ इसलिए चलती है कि $\sin(0\cdot\theta_k) = 0$, और पंक्ति $n$ इसलिए कि $\sin((n+1)\theta_k) = \sin(k\pi) = 0$: अर्थात्
सीमा-शर्तें ठीक $\theta_k = \frac{k\pi}{n+1}$ चुन लेती हैं। अतः $T_nv_k = 4\sin^2\bigl(\frac{k\pi}{2(n+1)}\bigr)v_k$; और $n$ मान $\intoo04$ में भिन्न हैं
तथा $v_k \neq 0$ भी: यही पूरा \omterm{def:b2:reduction:eigen}{वर्णक्रम} है, जो धनात्मक है, इसलिए $T_n$ धनात्मक
निश्चित है।

\textbf{22.} $\min_id_i\,I \leq D \leq \max_id_i\,I$, अतः $T_n +
\min_id_i\,I \leq T_n + D \leq T_n + \max_id_i\,I$ ($T_n$ जोड़ने से क्रम बना रहता है); और प्रश्न 5
तथा $\lambda_k(T_n + cI) = \lambda_k(T_n) + c$ वह घेरा दे देते हैं।

\textbf{23.} $T_2 = \begin{pmatrix} 2 & -1\\ -1 &
2\end{pmatrix}$ का \omterm{def:b2:reduction:eigen}{वर्णक्रम} $\{3, 1\}$ है, और
\[
2 - \sqrt2 \;\leq\; 1 \;\leq\; 2 \;\leq\; 3 \;\leq\; 2 +
\sqrt2 :
\]
अर्थात् कोशी अंतर्गुंथन सत्यापित। (ये वही \omterm{def:b2:reduction:eigen}{वर्णक्रम} हैं जो प्रश्न 15 में
थे: $\operatorname{diag}(1,-1,1)$ से संयुग्मन विकर्णेतर का चिह्न पलट देता है।) ग्राफ़ वाला पाठ:
$T_2$ उस पथ का लाप्लासीय-प्रकार का आव्यूह है जिसमें अंतिम शीर्ष हटा दिया गया
है --- अर्थात् कोई मुख्य उपआव्यूह, ठीक प्रश्न 11 वाली स्थिति।

\textbf{24.} चरम अभिलक्षणिक मानों का व्यवहार
\[
\lambda_{\min}(T_n) = 4\sin^2\frac{\pi}{2(n+1)}
\sim \frac{\pi^2}{(n+1)^2},
\qquad
\lambda_{\max}(T_n) = 4\cos^2\frac{\pi}{2(n+1)}
\longrightarrow 4 .
\]
है। अतः
\[
\kappa_n = \frac{\lambda_{\max}}{\lambda_{\min}}
\sim \frac{4(n+1)^2}{\pi^2} :
\]
यानी जाल जितना बारीक, विविक्त दूसरा अवकलज उतना ही दुरवस्थित --- और यही
तथ्य संख्यात्मक रैखिक बीजगणित की रूपरेखा तय करता है।

\textbf{25.} (क) किसी सामान्य आव्यूह के विक्षोभ से \omterm{def:b2:reduction:charpoly}{अभिलक्षणिक बहुपद} के
मूल उग्र रूप से हिल सकते हैं, परंतु न्यूनतम-महत्तम अभिलक्षण हर \omterm{def:b2:hermitian:adjoint}{हर्मीशियन}
\omterm{def:b2:reduction:eigen}{अभिलक्षणिक मान} को स्पष्ट अनुकूलन-मानों के बीच बाँध देता है, जिससे प्रश्न
8--9 की $1$-लिप्शिट्स स्थायित्व बाध्य हो जाती है। (ख) प्रश्न 1, 5, 16--20
में केवल चरम रैले मान लगे; जबकि वाइल, अंतर्गुंथन, शूर और काय फान (प्रश्न
7--14) को उपसमष्टियों पर पूरा न्यूनतम-महत्तम सचमुच चाहिए था। (ग) लोएव्नर
क्रम इन अदिश असमिकाओं को आव्यूह असमिकाओं के कलन में बदल देता है --- असली
जालों (प्रश्न 17) और असली प्रमेयों (प्रश्न 18--19) के साथ। (घ) ये औज़ार
संख्यात्मक विश्लेषण (प्रश्न 24 के दृढ़ता आव्यूह), क्वांटम विक्षोभ सिद्धांत
(वाइल: ऊर्जा-स्तर अधिकतम विक्षोभ के \omterm{def:b2:nvs:norm}{मानदंड} जितना ही खिसकते हैं), और तृतीय
वर्ष के खंड के \omterm{def:b2:metric:compact}{संहत} स्वसंलग्न संकारकों के न्यूनतम-महत्तम सिद्धांत की रोज़
की रोटी हैं।
\end{solution}
